\documentclass[10pt,a4paper]{amsart}
\usepackage[margin=19mm]{geometry}
\usepackage{amsmath,amssymb,mathtools}
\usepackage[scr=rsfs,cal=euler]{mathalfa}
\usepackage{xfrac}
\usepackage[unicode,hidelinks,bookmarks=false]{hyperref}
\newcommand{\ie}{i.e.\ }
\newcommand{\cf}{cf.\ }
\newcommand{\resp}{respectively\ }
\newcommand{\Subsection}{\subsection}
\providecommand{\nobreakdash}{\nobreak\hbox{-}}
\setlength{\emergencystretch}{2em}
\multlinegap0pt
\usepackage{bm}
\usepackage{bbm}
\usepackage{dsfont}
\usepackage{xspace}
\usepackage{cancel}
\usepackage{mathtools}
\usepackage{amsthm}
\makeatletter
\@ifundefined{theorem}{\newtheorem{theorem}{Theorem}}{}
\@ifundefined{lemma}{\newtheorem{lemma}[theorem]{Lemma}}{}
\makeatother
\title[Sherman-Takeda type theorems for locally C*-algebras]{Sherman-Takeda type theorems for locally C*-algebras}
\author{Lav Kumar Singh, Aljo\v{s}a Peperko}
\date{Source: arXiv:2601.00717v2 (CC BY 4.0)}
\begin{document}
\maketitle

\noindent\emph{Source and licence.} Sherman-Takeda type theorems for locally C*-algebras. Version arXiv:2601.00717v2 (CC BY 4.0), distributed under the Creative Commons Attribution 4.0 International licence, as recorded in the arXiv licence metadata for this exact version and shown on the abstract page https://arxiv.org/abs/2601.00717v2. Full rights notice: licenses/arxiv-2601.00717-CC-BY-4.0.txt.\par\smallskip

\noindent\textit{Editorial scope.} Counted unit: the Lemma whose source label is \texttt{biltolin} in the article (source0.tex lines 317--319, source label \texttt{biltolin}), delivered together with the article's own complete original proof of it (source0.tex lines 321--367, one \texttt{proof} environment of 47 lines). No mathematics is written, repaired or re-derived: statement and proof are the article's own text line for line. Required context retained uncounted: none. Every definition, display and label of those lines is retained unchanged. Omitted from this package and not redistributed: 1--316, 320--320, 368--469. Nothing inside a retained excerpt is omitted or altered: every retained body is delivered exactly as the article's lines are written, so no omission receipt of any kind accompanies this package. Citations: the retained excerpts carry no citation command at all, so no bibliography entry of the article is transcribed and no reference is redistributed. Reference adaptation: none was needed and none was made; every reference inside the delivered unit resolves inside it, so no unresolved reference or citation is produced. Printed slips kept literal: no printed word, symbol, hypothesis or number of the counted statement or of its proof was altered. Layout and class: the article's own class is \texttt{article} with packages including geometry, fullpage, amssymb, amsmath, hyperref, amsthm, amsfonts, graphicx, mathabx, cleveref, lipsum, authblk, enumitem, xcolor; the wrapper is the lane's pinned \texttt{amsart} editorial wrapper at 10pt on A4, which restates the theorem-like environments the retained text needs and adds the article's own macro definitions that the retained text needs (none), with extra wrapper packages (bm, bbm, dsfont, xspace, cancel, mathtools). The delivered numbering is the unit's own. Assets: no retained span contains a figure environment, a graphics-inclusion command, a PDF-page inclusion, a table or a TikZ picture; no image file is copied into this package, the package declares no asset dependency and no image file is redistributed with it. Licence: version arXiv:2601.00717v2 of this article is distributed under the Creative Commons Attribution 4.0 International licence, as recorded in the arXiv licence metadata for this exact version and shown on the abstract page https://arxiv.org/abs/2601.00717v2. A full rights notice is delivered at licenses/arxiv-2601.00717-CC-BY-4.0.txt. Characters: every retained excerpt is pure ASCII, so the delivered engine, XeLaTeX with its default Latin Modern text font, reports no missing character.
	\begin{lemma}\label{biltolin}
		Let $\mathcal U$ and $\mathcal V$ be locally convex topological vector spaces and $m:\mathcal U\times \mathcal V\to \mathcal V$ be a jointly continuous bilinear map. Then the induced map $\hat m:\mathcal U\to B(\mathcal V)$ is continuous, where $B(\mathcal V)$ is the space of continuous linear operators on $\mathcal V$ and it is equipped with topology of uniform convergence on bounded sets of $\mathcal V$.
	\end{lemma}

	\begin{proof}
		Since $\widehat{m}$ is linear (by bilinearity of $m$), it suffices to prove continuity at $0\in\mathcal{U}$. 
		We need to show that if $(u_\lambda)$ is a net in $\mathcal{U}$ with $u_\lambda \to 0$, then 
		$\widehat{m}(u_\lambda)\to 0$ in $B(\mathcal V)$. 			
		A basic neighbourhood of $0$ in $B(\mathcal V)$ is of the form
		\[
		\mathcal{W}(K,W)=\bigl\{T\in B(\mathcal V):T(K)\subseteq W\bigr\},
		\]
		where $K\subseteq \mathcal{V}$ is a bounded set and $W\subseteq \mathcal{V}$ is a convex balanced neighbourhood of $0$. 
		(We may assume without loss of generality that $K$ is absolutely convex.)		
		Fix such a neighbourhood $\mathcal{W}(K,W)$. We need to find $\lambda_0$ such that 
		$\widehat{m}(u_\lambda)\in\mathcal{W}(K,W)$ for all $\lambda\ge\lambda_0$, 
		i.e.,\ $m(u_\lambda,K)\subseteq W$.	Since $m$ is jointly continuous at $(0,0)$ and $m(0,0)=0$, there exist convex balanced 
		neighbourhoods $P\subseteq \mathcal{U}$ of $0$ and $Q\subseteq \mathcal{V}$ of $0$ such that
		\[
		m(P\times Q)\subseteq W.
		\]
		
		
		Since $K$ is bounded and $Q$ is a neighbourhood of $0$ in $\mathcal{V}$, there exists $t>0$ such that
		\[
		K\subseteq tQ.
		\]
		
		
		Define the scaled neighbourhood $			P':=\frac{1}{t}P=\bigl\{u\in\mathcal{U}:tu\in P\bigr\}.$
		Then $P'$ is a convex balanced neighbourhood of $0$ in $\mathcal{U}$.	Let $u\in P'$ and $v\in K$.  
		Then $tu\in P$ (write $tu=u'$ with $u'\in P$) and $v=tq$ for some $q\in Q$. 
		By bilinearity of $m$,
		\begin{align*}
			m(u,v)
			&=m\!\left(\frac{1}{t}u',tq\right)
			=\frac{1}{t}\,m(u',tq)
			=\frac{1}{t}\cdot t\cdot m(u',q)
			=m(u',q).
		\end{align*}
		Since $u'\in P$ and $q\in Q$, we have $m(u',q)\in W$, hence $m(u,v)\in W$. 
		As $v\in K$ was arbitrary, $m(u,K)\subseteq W$.				
		Since $u_\lambda\to 0$ in $\mathcal{U}$, there exists $\lambda_0$ such that 
		$u_\lambda\in P'$ for all $\lambda\ge\lambda_0$. 
		For these indices, we have
		\[
		\widehat{m}(u_\lambda)(K)=m(u_\lambda,K)\subseteq W,
		\]
		so $\widehat{m}(u_\lambda)\in\mathcal{W}(K,W)$. 
		Thus $\widehat{m}(u_\lambda)\to 0$ in $B(\mathcal V)$.
	\end{proof}

\end{document}
